operations @@M0@@ define, by 7.3, additive endomorphisms of the functor @@M2@@ with values in @@M1@@-rings; by 1.10, they are characterized by their value on @@M3@@. We have:
@@M4@@
Consequently:
\sgaStableTarget{sga6.e05.eq.7-4-1}@@M5@@

\sgaStableTarget{sga6.e05.prop.7-5}\textbf{Proposition 7.5.} i) \uline{If} @@M6@@ \uline{is a} @@M7@@-\uline{ring}, \uline{the operations} @@M8@@ \uline{satisfy} \uline{the identities}:
\sgaStableTarget{sga6.e05.eq.7-5-1}@@M9@@
ii) \uline{Conversely}, \uline{if the} @@M10@@ \uline{satisfy these identities}, \uline{and if} @@M11@@ \uline{is of} \uline{characteristic zero}, @@M12@@ \uline{is a} @@M13@@-\uline{ring}.

i) If @@M14@@ is a @@M15@@-ring, we have @@M16@@, whence: @@M17@@, and consequently:
@@M18@@

To prove the second of relations (7.5.1), we consider the diagram:
\sgaStableTarget{sga6.e05.eq.7-5-2}@@M19@@
where the vertical arrows are defined by multiplication, and the maps @@M20@@ as follows: the ``logarithmic derivative'' map is
